% Part: sets-functions-relations
% Chapter: induction
% Section: induction

\documentclass[../../../include/open-logic-section]{subfiles}

\begin{document}

\olfileid{sfr}{ind}{int}

\olsection{مقدمه}

\begin{explain}
استقرا فن برهانی پرکاربرد در ریاضیات و منطق است. شاید آن را از ریاضیات
بشناسید، جایی که بر اعداد طبیعی استقرا می‌کنیم. در واقع، گونه‌های فراوانی
از مجموعه‌های اشیای ریاضی و منطقی وجود دارند که برای استقرا مناسب‌اند.

به‌طور کلی، استقرا امکان می‌دهد ادعایی کلی را اثبات کنیم؛ یعنی نشان دهیم
هر شیء از نوعی معین ویژگی‌ای دارد. استقرا به‌ویژه هنگامی سودمند است که
روش برهانِ «معرفی کلی» بیش از حد پیچیده باشد. استقرا تنها بر اشیای ریاضی‌ای
کار می‌کند که به شیوه‌هایی ویژه ساخته شده‌اند: اگر هر عضو مجموعه یا پایه
باشد یا از اعضای پایه ساخته شده باشد، می‌توان بر آن استقرا کرد. دقیق‌تر:
\end{explain}

\begin{defn}
مجموعهٔ~$S$ تحت تابع~$n$-موضعی~$f$ \emph{بسته} است اگر و تنها اگر
هرگاه~$x_1,\dots,x_n\in S$، آنگاه
$f(x_1,\dots,x_n)\in S$.
\end{defn}

\begin{explain}
برای نمونه، اعداد طبیعی (~$\mathbb{N}$) تحت تابع جانشین
$f(x) = x+1$ بسته‌اند. اگر~$x$ عددی طبیعی باشد، $x+1$ نیز طبیعی است.
اما اعداد طبیعی تحت تابع پیشین~$f(x) = x-1$ بسته نیستند، زیرا~$0$ عددی
طبیعی است ولی~$-1$ طبیعی نیست.
\end{explain}

\begin{defn}
می‌گوییم ویژگی~$P$ تحت تابع~$n$-موضعی~$f$ \emph{حفظ می‌شود} هرگاه برای
هر~$a_1,\dots,a_n$ در دامنهٔ~$f$، از
$P(a_1),\dots,P(a_n)$ نتیجه شود
$P(f(a_1,\dots,a_n))$؛ یعنی اگر همهٔ ورودی‌ها ویژگی~$P$ را داشته باشند،
خروجی~$f$ نیز آن را داشته باشد.
\end{defn}

\begin{explain}
ویژگی «زوج‌بودن» تحت تابع~$f(x) = x+2$ حفظ می‌شود، زیرا اگر~$x$ زوج
باشد، $x+2$ نیز زوج است. ویژگی زوج‌بودن تحت تابع جانشین حفظ نمی‌شود،
زیرا اگر~$x$ زوج باشد، $x+1$ فرد است.

استقرا بر اعداد طبیعی کار می‌کند، زیرا با تعداد متناهی کاربردِ تابع جانشین
می‌توان از~$0$ به هر عدد طبیعی رسید. این ویژگی هر مجموعه‌ای است که بتوان
بر آن استقرا کرد.
\end{explain}

\begin{defn}[استقرا بر~$\Nat$]
اگر~$0$ ویژگی~$P$ را داشته باشد و~$P$ تحت تابع جانشین حفظ شود، آنگاه
هر عدد طبیعی ویژگی~$P$ را دارد.
\end{defn}

\begin{ex} مجموع اعداد طبیعی از~$0$ تا~$n$ برابر
$\frac{n(n+1)}{2}$ است. \\

\textbf{حالت پایه.} مجموع اعداد طبیعی از~$0$ تا~$0$ برابر
$0=\frac{0\cdot 1}{2}$ است.\\

\textbf{گام استقرا.} فرض کنید ویژگی برای~$k$ برقرار باشد. نشان می‌دهیم
برای~$k+1$ نیز برقرار است. به بیان دیگر، فرض کنید~$P(k)$ برقرار باشد؛
یعنی
\[ 0 + 1 + \ldots + k = \frac{k(k+1)}{2} \]
با استفاده از فرض~$P(k)$ نشان می‌دهیم~$P(k+1)$ برقرار است؛ یعنی
\[ 0 + 1 + \ldots + k + (k+1) = \frac{(k+1)(k+2)}{2}. \]

\begin{align*}
0 + 1 + \ldots + k &= \frac{k(k+1)}{2} \tag{فرض} \\
0 + 1 + \ldots + k + (k+1) &= \frac{k(k+1)}{2} + (k+1)
\tag{افزودن $k+1$ به دو طرف} \\
&= \frac{k(k+1) + 2(k+1)}{2} \\
&= \frac{k^2 + k + 2k +2}{2} \\
&=\frac{(k+1)(k+2)}{2}
\end{align*}
پس گام استقرا برقرار است و ادعا را اثبات کرده‌ایم.
\end{ex}

\begin{explain}
برای~!!{formula}s، اعضای پایه همان~!!{formula}s اتمی‌اند. !!{formula}s
پیچیده‌تر با پیوندزدن~!!{formula}s ساده‌تر به‌وسیلهٔ عملگرها ساخته
می‌شوند. هر عملگر را می‌توان متناظر با تابعی از~!!{formula}s دانست که
!!{formula} تازه‌ای بازمی‌گرداند و ورودی‌ها را با عملگر مورد نظر به هم
می‌پیوندد. برای نمونه، تابعی که دو~!!{formula}s~$!A$ و~$!B$ را به صورت
عطف به هم می‌پیوندد، چنین نوشته می‌شود:
$\mathcal E _\land (!A, !B) = !A \land !B$. این‌ها را
\emph{توابع فرمول‌ساز} می‌نامیم. مجموعهٔ~!!{formula}s تحت این توابع بسته
است: هرگاه~$!A$ و~$!B$، !!{formula}s باشند، $\lnot !A$،
$!A \land !B$ و مانند آن‌ها نیز از همین مجموعه‌اند.
\end{explain}

\begin{defn}[استقرا بر~!!{formula}s]
اگر هر~!!{formula} اتمی ویژگی~$P$ را داشته باشد و~$P$ تحت توابع
!!{formula}-ساز حفظ شود، آنگاه هر~!!{formula} ویژگی~$P$ را دارد.
\end{defn}

\begin{explain}
این تعریف «دستور کاری» برای برهان‌های استقرایی بر~!!{formula}s پیشنهاد
می‌کند. اگر خواسته شود نشان دهیم هر~!!{formula} ویژگی~$P$ را دارد، مراحل
زیر را دنبال می‌کنیم:\\

\textbf{حالت پایه.} فرض کنید~$!A$ یک~!!{formula} اتمی باشد. [\ldots]
پس~$!A$ ویژگی~$P$ را دارد.

\textbf{گام استقرا.} فرض کنید~$!A$ و~$!B$ دو~!!{formula}s باشند که هر
دو ویژگی~$P$ را دارند.

حالت~$\lnot$: [\ldots] پس~$\lnot !A$ ویژگی~$P$ را دارد.

حالت~$\land$: [\ldots] پس~$!A \land !B$ ویژگی~$P$ را دارد.

حالت~$\lor$: [\ldots] پس~$!A \lor !B$ ویژگی~$P$ را دارد.

حالت~$\lif$: [\ldots] پس~$!A \lif !B$ ویژگی~$P$ را دارد.

حالت~$\liff$: [\ldots] پس~$!A \liff !B$ ویژگی~$P$ را دارد.

حالت~$\lforall[]$: [\ldots] پس~$\lforall[x] !A$ ویژگی~$P$ را دارد.

حالت~$\lexists[]$: [\ldots] پس~$\lexists[x] !A$ ویژگی~$P$ را دارد. \\

بنابراین هر~!!{formula} ویژگی~$P$ را دارد.
\end{explain}

\end{document}
